% Appendix A:
%\subsection{Introduction}
%\subsection{Model of an Abelian Scheme}

Our goal in this appendix is to present a treatment of the
Silverman--Tate Theorem, \cite[Theorem~A]{Silverman}, using the language
of Cartier divisors. Using Cartier divisors as opposed to Weil
divisors  allows us to relax the flatness 
hypotheses imposed on $\pi$ in the notation of \cite[$\mathsection$3]{Silverman}. Apart from this minor tweak
we closely follow the original argument presented by Silverman. 

Suppose $S$ is a {regular},
{irreducible}, {quasi-projective} variety over $\IQbar$.
Let $\pi \colon \cA\rightarrow S$ be an abelian scheme.
We write $\eta$ for the generic point of $S$ and $\cA_\eta$ for the
generic fiber of $\pi$. Then $\cA_\eta$ is an abelian variety defined
over $\IQbar(\eta)$. 

Suppose we are presented with a closed immersion
$\cA\rightarrow \IP_{\IQbar}^n\times S$ over $S$ and
 with a projective variety $\overline S$
containing $S$ as a Zariski open and dense subset.
We will assume that $\overline S$ is embedded into $\IP_\IQbar^m$.
We do not assume that $\overline S$ is regular. % Let $\cM$ be a fixed
% ample line bundle on $\overline S$.

We identify $\cA$ with a subvariety of $\IP_\IQbar^n\times S$.
Moreover, let $\overline{\cA}$ denote the Zariski closure of 
$\cA$ in $\IP_S^n=\IP_{\IQbar}^n\times\overline S \subset
\IP_{\IQbar}^n\times\IP_\IQbar^m$. %% We consider $S$
%% as Zariski open in $\overline S$ and $\cA$ as Zariski open in
%% $\overline \cA$.

We set $\overline\cL = \cO(1,1)|_{\overline{\cA}}$ and  $\cL
= \overline\cL|_{\cA}$. We  will  assume in addition that 
$[-1]^*\cL_\eta \cong \cL_\eta$ 
where  $\cL_\eta$ is the restriction of $\cL$ to $\cA_\eta$. 
This implies
$[2]^*\cL_\eta \cong \cL_\eta^{\otimes 4}$.

\newcommand{\hcA}{h}

Given these immersions, we have several height functions.
For
$(P,s)\in \overline{\cA}(\IQbar)\subset \IP^n_\IQbar(\IQbar)\times \IP^m_\IQbar(\IQbar)$ we define
$\hcA(P,s) = h(P)+h(s)$ using the Weil height.
Moreover, for $s\in \overline S(\IQbar)\subset \IP^m_\IQbar(\IQbar)$
we define $h_{\overline S}(s) = h(s)$. 
Finally, 
for all  $P\in \cA(\IQbar)$ we denote by
  \begin{equation*}
    \hat h_{\cA}(P) = \lim_{N\rightarrow \infty}
    \frac{\hcA([N](P))}{N^2}
  \end{equation*}
the N\'eron--Tate height with respect to $\cL$;  it is
well-known that the limit converges, \textit{cf.} the reference around
(\ref{eq:fiberwiseNT}). 
  
We will prove the following variant of the Silverman--Tate Theorem.

\begin{theorem}
  \label{thm:silvermantate}
  There exists a constant $c>0$  such that for all 
  $P \in \cA(\IQbar)$ we have 
  \begin{equation*}
    \bigl| \hat h_{\cA}(P) -
    \hcA(P)\bigr| \le c  \max\{1,h_{\overline
      S}(\pi(P))\}. 
  \end{equation*}
\end{theorem}

The constant $c$ depends on $\cA$ and on the various immersions but not on $P$. 
The proof is distributed over  the  next subsections.

\subsection{Extending multiplication-by-$2$}
We keep the notation from the  previous subsection.
We have  constructed a (very naive) projective model $\overline\cA$ of $\cA$.
Note that $\overline\cA$ and $\overline S$ may fail to be regular.
Moreover, the natural morphism
$\overline\cA\rightarrow \overline S$, which we also denote by $\pi$,  may fail to be smooth or even
flat.

% Let us first desingularize the compactified base $\overline S$ by applying
% Hironaka's Theorem. Thus there is a proper, birational morphism
% $b\colon \overline S'\rightarrow\overline S$ that is an isomorphism
% above $S$
% and such that  $\overline S'$
% is regular. Note that $b$ is even
% projective and $\overline S'$ is integral. So $\overline S'$ is an
% irreducible, regular, projective variety. 

% Now consider the base change $\overline
% \cA\times_{\overline S}\overline S'$. This new scheme
% may not be a suitable model as it may fail to be irreducible or even
% reduced. However, recall that $b$ is an isomorphism above $S\subset
% \overline S$. So $(\overline
% \cA\times_{\overline S}\overline S')_S =
% %\overline{\cA}'\times_{\overline{S}'}S=
% \overline{\cA}\times_{\overline S}S$  is isomorphic to $\cA$ and thus
% integral. We may consider $\cA$ as an open subscheme of
% $\overline
% \cA\times_{\overline S}\overline S'$. It must be contained in an irreducible component of
% $\overline
% \cA\times_{\overline S}\overline S'$. We endow this irreducible component with the
% reduced induced structure and obtain an integral, closed subscheme
% $\overline{\overline{\cA}}\subset \overline
% \cA\times_{\overline S}\overline S'$. Its image in
% $\overline{S}'$ contains the open subset $b^{-1}(S)$. As $\overline S'$
% is irreducible, we find by properness that
% $\overline{\overline{\cA}}\rightarrow\overline S'$ is surjective, it is
% also projective. 

% To simplify notation we forget about $\overline{\cA}$ and denote
% $\overline{\overline{\cA}}$ by $\overline{\cA}$. 

% Let us recapitulate:  we are presented with projective varieties
% $\overline{\cA}$ and $\overline{S}$ containing $\cA$ and $S$ as
% Zariski open subsets, respectively, such that $\pi$ extends to a morphism 
%  $\overline{\cA}\rightarrow \overline
% S$.
% Moreover, $\overline{S}$ is regular. 

% However,    $\overline\cA$  may fail to be 
% regular and 
% $\overline\cA\rightarrow \overline S$  may fail to be flat. 

% Let us work on $\overline \cA$ to solve the first problem.

Multiplication-by-$2$ is a morphism $[2]\colon \cA\rightarrow \cA$
that extends to a rational map
$\overline{\cA}\dashrightarrow\overline{\cA}$.
We consider the graph of $[2]$ on $\cA$ as a subvariety of
$\cA\times_S\cA$.
Let $\overline{\cA}'$ be the Zariski closure of this graph inside
$\overline{\cA}\times_{\overline S}\overline{\cA}$. 
 Write $\rho\colon
\overline{\cA}'\rightarrow\overline{\cA}$ for the restriction of the projection onto the
first factor and  $[2]$ for the restriction onto the second factor.
We may identify $\cA$ with a Zariski open subset of
$\overline{\cA}'$. Under this identification, $\rho$ restricts to the
identity on $\cA$ and $[2]$ restricts to multiplication-by-$2$ on
$\cA$. 
%% There is a blowup
%% $\overline{\cA}' \rightarrow \overline{\cA}$ with center outside $\cA$
%% such that $[2]$ extends to a morphism $[2] \colon \overline{\cA}'
%% \rightarrow \overline{\cA}$.

% By composing we get a morphism $[2] \colon
% \widetilde{\cA} \rightarrow \overline {\cA}'$ whose restriction to $\cA$ is
% multiplication-by-$2$.
The following diagram  commutes
\begin{equation*}
  \xymatrix{
%&    &  \widetilde{\cA} \ar[d]_{\beta} & \\
 & \overline{\cA} \ar@{}[d]|-*[@]{\supset}      &\ar[l]_\rho \overline{\cA}' \ar@{}[d]|-*[@]{\supset}   \ar[r]^{[2]} & \overline{\cA}\ar@{}[d]|-*[@]{\supset} \\
& \cA \ar@{=}[r]  & {\cA}      \ar[r]^{[2]} &{\cA} 
  }  
\end{equation*}
where the first and third inclusions are equal and the middle one
comes from the identification involved in the graph construction. 



\subsection{Proof of the Silverman--Tate Theorem}


% Let $K=\IQbar$  and suppose $S$ is a regular, 
% irreducible, quasi-projective variety over $K$. Let $\pi\colon \cA\rightarrow S$
% be an abelian scheme.
We keep the notation from the  previous subsection.


\begin{proposition}
  There exists a constant $c_1>0$ such that 
  \begin{equation}
    \label{eq:duplicationbound}
\left|   \hcA([2](P)) - 4
   \hcA(P)\right|  \le c_1  \max\{1,h_{\overline
   S}(\pi(P))\}
  \end{equation}
  holds for all $P\in \cA(\IQbar)$. 
\end{proposition}
\begin{proof}
  We define
  \begin{equation}
  \label{def:cF2}
  \cF' =  [2]^* \overline{\cL} \otimes
  \rho^*\overline{\cL}^{\otimes(-4)} \in \mathrm{Pic}(\overline{\cA}'). 
\end{equation}
  
Recall that we have identified $\cA$ with a Zariski open subset of
$\overline{\cA}'$. 
The  restriction of $[2]^*\overline{\cL}$ to the generic fiber
$\cA_\eta\subset\cA\subset\overline{\cA}'$
coincides with $[2]^*\cL_\eta$ and the restriction of $\rho^*\overline{\cL}$  to $\cA_\eta$
is identified with $\cL_\eta$. Using our assumption
 $[2]^*\cL_\eta \cong
\cL_\eta^{\otimes 4}$ on the generic fiber $\cA_\eta$ we see that  $\cF'$ is trivial on $\cA_\eta$.

By \cite[Corollaire~21.4.13 (pp.~361 of EGA IV-4, in Errata et Addenda, liste~3)]{EGAIV} %% (Errata et appenda)
applied to
$\cA\rightarrow S$ there exists a line
bundle $\cM$ on $S$ such that $\pi|_{\cA}^*\cM \cong \cF'|_{\cA}$.

Let us first desingularize the compactified base $\overline S$ by applying
Hironaka's Theorem. Thus there is a proper, birational morphism
$b\colon \overline S'\rightarrow\overline S$ that is an isomorphism
above $S$
 such that  $\overline S'$
is regular. We consider $S$ as Zariski open in $\overline S'$.  Note that $b$ is even
projective and $\overline S'$ is integral. So $\overline S'$ is an
irreducible, regular, projective variety. 

Now consider the base change $\overline
\cA'\times_{\overline S}\overline S'$. This new scheme
may  fail to be irreducible or even
reduced. However, recall that $b$ is an isomorphism above the regular $S\subset
\overline S$. So $(\overline
\cA'\times_{\overline S}\overline S')_S =
%\overline{\cA}'\times_{\overline{S}'}S=
\overline{\cA}'\times_{\overline S}S$  is isomorphic to $\cA$ and thus
integral. We may consider $\cA$ as an open subscheme of
$\overline
\cA'\times_{\overline S}\overline S'$. It must be contained in an irreducible component of
$\overline
\cA'\times_{\overline S}\overline S'$. We endow this irreducible component with the
reduced induced structure and obtain an integral, closed subscheme
$\overline{\overline{\cA}}\subset \overline
\cA'\times_{\overline S}\overline S'$. We get a commutative diagram
\begin{equation*}
  \xymatrix{
\cA\ar@{}[r]|-*[@]{\subset}\ar[d]_\pi &
 \overline{\overline{\cA}} \ar[r]^f \ar[d]_{\overline\pi} & \overline{\cA}' \ar[d]_{\pi\circ\rho}&\\
S \ar@{}[r]|-*[@]{\subset} &\overline{S}' \ar[r] &      \overline{S} &
  }  
\end{equation*}
The horizontal  morphisms compose to the identity on the domain. % We write $f\colon \overline{\overline\cA}\rightarrow \overline \cA'$ for
% the natural morphism. 
% Its image in
% $\overline{S}'$ contains the open subset $\beta_1^{-1}(S)$. As $\overline S'$
% is irreducible, we find by properness that
% $\overline\pi\colon\overline{\overline{\cA}}\rightarrow\overline S'$ is surjective, it is
% also projective. 

We consider $S$ as a Zariski open subset of $\overline S'$. 
As $\overline S'$ is regular, we can extend $\cM$ to a line bundle on
the regular $\overline S'$, \textit{cf.} \cite[Corollary~11.41]{GoertzWedhorn}. % G\"ortz--Wedhorn.
The pull-back $f^*\cF'\otimes\overline\pi^*\cM^{\otimes(-1)}$ is
trivial on $\cA\subset\overline{\overline{\cA}}$. 

By Hironaka's Theorem there is a proper, birational morphism $\beta\colon\widetilde
{\cA}\rightarrow\overline{\overline\cA}$ that is an isomorphism above $\cA$
(which is regular) 
such that $\widetilde{\cA}$ is regular. We may identify $\cA$ with a Zariski open subset of
$\widetilde \cA$.
%But it may not be trivial on $\widetilde \cA$. 

Now we pull everything back to the regular  $\widetilde\cA$.
More precisely, we set $\cF = \beta^*f^*\cF'$.
% We consider $\cA$ as a Zariski open subset of $\widetilde\cA$.
Then
$\cF\otimes\beta^*\overline\pi^*\cM^{\otimes(-1)}$ is trivial when restricted to
$\cA$. 

To a Cartier divisor $D$ we attach its line bundle $\cO(D)$.
As $\widetilde \cA$ is integral we may fix a Cartier divisor $D$ on
$\widetilde \cA$ with $\cO(D) \cong \cF\otimes \beta^*\overline\pi^*
\cM^{\otimes(-1)}$. Let $\mathrm{cyc}(D)$ denote the Weil divisor of $\widetilde\cA$
attached to $D$. The linear equivalence class of $\mathrm{cyc}(D)$
restricted to $\cA$ is trivial. By \cite[Proposition~11.40]{GoertzWedhorn}
%% \footnote{This proposition is false as stated, one
%%   needs \textit{e.g.} in addition that $X=\widetilde{\cA}$ is reduced. Recall
%%   that the  blowup of a reduced scheme is reduced.}
$\mathrm{cyc}(D)$ is linearly equivalent to a Weil divisor $\sum_{i=1}^r n_i
Z_i$ with  $Z_i \subset\widetilde{\cA}\ssm \cA$ irreducible and of
codimension $1$ in $\widetilde{\cA}$. 

We let $\widetilde\pi$ denote the composition
$\widetilde\cA\rightarrow \overline{\overline\cA}\rightarrow \overline{S}'$. 
Let us consider $\widetilde\pi(Z_i) = Y_i$. As $\widetilde\pi$ is proper, each $Y_i$ is
an irreducible closed subvariety of $\overline S'$. Moreover, $Y_i
\subset \widetilde\pi(\widetilde{\cA}\ssm \cA) \subset \overline S'\ssm S$. So $Y_i$ has
dimension at most $\dim \overline S' - 1$. But $Y_i$ could have
codimension at least $2$ and thus fail to be the support of a Weil divisor.
On the regular $\overline S'$ a 
Cartier divisor is the same thing as a Weil divisor;
see \cite[Theorem~11.38(2)]{GoertzWedhorn}.
For each $i$ we fix a Cartier divisor  $E_i$ of $\overline S'$
such that   $\mathrm{cyc}(E_i)$ equals a prime Weil divisor supported on an irreducible
subvariety containing $Y_i$.
 Since $\mathrm{cyc}(E_i)$ is effective, we find that $E_i$ is effective, see \cite[Theorem~11.38(1)]{GoertzWedhorn}
 and its proof. 
An effective Cartier divisor and its image under the cycle map
$\mathrm{cyc}(\cdot)$ have equal support. So the subscheme of $\overline S'$ attached to  $E_i$
contains $Y_i$. 

The pull-back $\widetilde\pi^*E_i$ is well-defined as a Cartier
divisor, we do not require that $\pi$ is flat, \textit{cf.}
\cite[Proposition~11.48(b)]{GoertzWedhorn}. 
By \cite[Corollary~11.49]{GoertzWedhorn} the inverse image
$\widetilde\pi^{-1}(E_i)$, taken as  a subscheme of $\widetilde{\cA}$
is the subscheme attached to  $\widetilde\pi^*E_i$ and $\widetilde\pi^*E_i$ is
effective. 

Note that $\widetilde\pi^{-1}(E_i)\supset \widetilde\pi^{-1}(Y_i)\supset Z_i$. 
The support satisfies $\mathrm{Supp}(\widetilde\pi^* E_i)\supset Z_i$.
Moreover, as $\widetilde\pi^*E_i$ is effective,  $\mathrm{cyc}(\widetilde\pi^* E_i)$ is effective and
 $\mathrm{Supp} (\mathrm{cyc}(\widetilde\pi^* E_i)) =
\mathrm{Supp}(\widetilde\pi^*E_i)$.
%% The latter requires Exercise 11.16 in GW, no regularity
%% condition on \widetilde\cA is needed.
Thus
%$\mathrm{cyc}(\pi^* E_i) \ge Z_i$. Then
\begin{equation*}
\pm   \sum_{i=1}^r n_i Z_i \le \mathrm{cyc}\left(\widetilde\pi^*\sum_{i=1}^r |n_i|
  E_i\right).
\end{equation*}

Recall that $\mathrm{cyc}(D) = \mathrm{cyc}(\mathrm{div}\phi) +
\sum_{i=1}^r n_i Z_i$ for some rational function $\phi$ on
$\widetilde{\cA}$. Therefore,
%% \begin{equation*}
%%   \pm \mathrm{cyc}(D - \mathrm{div}(f)) \le
%%   \mathrm{cyc}\left(\widetilde\pi^*\sum_{i=1}^r |n_i| E_i\right).
%% \end{equation*} 
% Thus the support of $\mathrm{cyc}(D-\mathrm{div}(f))$ is contained in
% the support of $  \mathrm{cyc}\left(\pi^*\sum_{i=1}^r |n_i|
%   E_i\right)$. 
%% Thus
\begin{equation*}
%  \pm \mathrm{cyc}(D - \mathrm{div}(f)) \le
0\le   \mathrm{cyc}\left(\pm (D-\mathrm{div}\phi) + \widetilde\pi^*\sum_{i=1}^r |n_i| E_i\right).
\end{equation*}

Since $\widetilde{\cA}$ is regular and in particular normal, we find  that 
\begin{equation}
\label{eq:thisiseffective}\pm(D - \mathrm{div}\phi) + \widetilde\pi^*\sum_{i=1}^r |n_i|E_i
\end{equation}
is an effective Cartier divisor for both  signs;
see \cite[Theorem~11.38(1)]{GoertzWedhorn} and its proof.
%% \footnote{I think this is
%%   the only point where we require some regularity property of
%%   $\widetilde\cA$. So perhaps it suffices to normalize $\overline{\overline\cA}$.}
Moreover, its support equals the support of 
$$0\le \mathrm{cyc}\left(\pm (D-\mathrm{div}\phi) + \widetilde\pi^*\sum_{i=1}^r |n_i|
  E_i\right) =
\pm \mathrm{cyc}(D-\mathrm{div}\phi) + \sum_{i=1}^r |n_i|\mathrm{cyc}(\widetilde\pi^*E_i).$$
Thus the support of (\ref{eq:thisiseffective}) lies in
$\bigcup_{i=1}^r \mathrm{Supp}(\widetilde\pi^* E_i)$. 


We apply $\cO(\cdot)$ and pass  again  to line bundles.
Let us denote $\mathcal{E} = \mathcal{O}(\sum_{i=1}^r |n_i| E_i)$, a
line bundle on $\overline S'$. The line bundle attached to 
 (\ref{eq:thisiseffective}) is $(\mathcal{F}\otimes\beta^*\overline\pi^*\mathcal{M}^{\otimes(-1)} 
 )^{\otimes(\pm 1)}\otimes\widetilde\pi^*\mathcal{E}$.
 Since (\ref{eq:thisiseffective}) is effective, both
$(\mathcal{F}\otimes\beta^*\overline\pi^*\mathcal{M}^{\otimes(-1)} 
 )^{\otimes(\pm 1)}\otimes \widetilde\pi^*\mathcal{E}$
% $\pi^*(\cE\otimes\cM^{\otimes (\mp 1)})\otimes\cF^{\otimes(\pm 1)}$ 
have a non-zero global
 section.

By the Height Machine this translates to 
 \begin{equation*}
   h_{\widetilde \cA,(\mathcal{F}\otimes\beta^*\overline\pi^*\mathcal{M}^{\otimes(-1)} 
 )^{\otimes(\pm 1)}\otimes \widetilde\pi^*\mathcal{E}}(P)\ge O(1)
 \end{equation*}
 for all $\widetilde P\in \widetilde\cA(\IQbar)$ 
with $\widetilde\pi(\widetilde P)\not\in\bigcup_i \mathrm{Supp}(E_i)$. 
% where $U = \widetilde\cA\ssm \bigcup_i
 % \mathrm{Supp}(\pi^* E_i)$. 
By functoriality properties of  the Height Machine we obtain 
 \begin{equation*}
  |h_{\overline \cA',\cF'}(f(\beta(\widetilde P)))| \le h_{\overline S',\cE}(\widetilde\pi(\widetilde P)) + |h_{\overline
    S', \cM}(\overline{\pi}(\beta(\widetilde P)))|
  +O(1)
 \end{equation*}
 for the same $\widetilde P$. 
 We recall (\ref{def:cF2}) and again use the Height Machine to find
 \begin{equation*}
\left|   \hcA([2](P')) - 4
   \hcA(\rho(P'))\right|   \le h_{\overline
   S',\cE}(\widetilde\pi(\widetilde P)) + |h_{\overline S',
     \cM}(\widetilde\pi(\widetilde P))|
  +O(1)
 \end{equation*}
where $P'=f(\beta(\widetilde P))$. Observe that all points of $\cA(\IQbar)$ are
in the image of $f\circ \beta$. 

We recall that the desingularization morphism $\overline S'\rightarrow\overline S$ is an
isomorphism above $S$ and that we have identified $\cA$ with a Zariski
open subset of $\overline{\cA}'$ and of $\overline{\cA}$.
Under these identifications and if $P'$ corresponds to $P\in
\cA(\IQbar)$, then $[2](P')$ is the duplicate of $P$,
$\rho(P')= P$, and $\widetilde\pi(\widetilde P) = \pi(\rho(P'))=\pi(P)$.  We apply the Height
 Machine a final time and use that $h_{\overline S}$ arises from the
 Weil height restricted to $\overline{S}(\IQbar)$. We
 find 
 \begin{equation*}
%\label{eq:preliminary1}
\left|   \hcA([2](P)) - 4
   \hcA(P)\right|   \le c_1
 \max\{1,h_{\overline S}(\pi(P))\}
 \end{equation*}
 for all $P\in \cA(\IQbar)$   with  $\widetilde\pi(\widetilde
 P)\not\in \bigcup_i \mathrm{Supp}(E_i)$, under the identifications
 above. 

Let $P\in\cA(\IQbar)$. As the $Y_i$ lie in $\widetilde\pi(\widetilde\cA\ssm
\cA)$ we can choose all $E_i$ above to avoid $\pi(P)$.  After doing this finitely often (using
noetherian induction) and replacing the $E_i$ from before and
adjusting $c_1$, we find 
\begin{equation*}
\left|   \hcA([2](P)) - 4
   \hcA(P)\right|  \le c_1 \max\{1,h_{\overline S}(\pi(P))\}
\end{equation*}
for all $P\in \cA(\IQbar)$ where $c_1>0$ is independent of $P$.
\end{proof}

\begin{proof}[Proof of Theorem~\ref{thm:silvermantate}]
  Having (\ref{eq:duplicationbound}) at our disposal the proof follows
  a well-known argument.   
  Indeed, say $l\ge k\ge 0$ are integers.
  Then applying the triangle inequality to the appropriate telescoping
  sum yields
  \begin{alignat*}1
    \left|\frac{\hcA([2^l](P))}{4^l}- \frac{\hcA([2^k](P))}{4^k}\right|
    &\le \sum_{m=k}^{l-1}
        \left|\frac{\hcA([2^{m+1}](P))}{4^{m+1}}-
        \frac{\hcA([2^m](P))}{4^m}\right|\\
        &\le \sum_{m=k}^{l-1}
        4^{-(m+1)}    \left|{\hcA([2^{m+1}](P))}- 4{\hcA([2^m](P))}\right|.
  \end{alignat*}
  We apply (\ref{eq:duplicationbound}) to $[2^m](P)$ and find that the
  sum is bounded by $c_1 x \sum_{m=k}^{l-1}4^{-(m+1)} \le c_1 x 4^{-k}$
  where $x = \max\{1,h_{\overline S}(\pi(P))\}$.
  So $\left(\hcA([2^l](P))/4^l\right)_{l\ge 1}$ is a Cauchy sequence
  with limit
  $\hat h_{\overline\cA}(P)$. 
  Taking $k=0$ and $l\rightarrow \infty$ we obtain from the estimates
  above that
  $|\hat h_{\overline\cA}(P) - \hcA(P)|\le c_1x$, as
  desired.   
\end{proof}
% \bibliographystyle{alpha}
% \bibliography{literature}

% \end{document}

%%% Local Variables:
%%% TeX-master: "main"
%%% End:


